\subsection{局所化で一段消える原点根基の厚み}

$r=1$ の通常CY族では、原点根基 $\mathcal R=(\Delta)/\mathfrak m_0\mathcal S$ の
正確な冪零指数は $j+4=(m+7)/2$ である。
$x$ 局所化後の $j+3=(m+5)/2$ は全環の指数ではない。
再利用する機構は、星積イデアルをPBW多項式イデアルと同定してから、
上界を明示的な所属式、下界を可換多項式の評価で分離することである。

$f=xyz-2y^3/3$、$g=z^2(4xz-3y^2)/9$ とすると
\[
 \Delta^2-x(xw^3+q^jyz w^2)J=q^2T_j
\]
であり、$T_j$ は公表稿XIの明示式で与えられる多項式である。
これに $\Delta^{j+2}$ を掛けて上界を得る。
$y=z=0$ では原点イデアルが $(x^{2j+7}w^{2j+5})$ へ写り、
$\Delta^{j+3}$ の像がそのイデアルに入らないので下界を得る。
評価を非可換環の準同型として用いない点を維持する。

この指数の問題は解決済みである。
次に調べるなら、根基の全ての層 $\mathcal R^a/\mathcal R^{a+1}$ の
境界上の台と加群構造、または動く中心指標との貼合せを対象とする。
